\chapter{Evolución Arquitectónica}
\label{cap:EvolucionArquitectonica}

\section{Exploración de Paradigmas: Clasificación vs. Regresión}
\subsection{Enfoque de Clasificación}
La clasificación es el paso previo natural a la búsqueda de un modelo de inferencia de parámetros.
De esta manera, se puede comprobar de forma práctica que la red neuronal procesa el audio adecuadamente
y es capaz de conservar su información acústica fundamental.

El \textbf{Prototipo 1} es una prueba de concepto para la primera aproximación al aprendizaje automático aplicado al audio.
En él se utiliza un dataset reducido, construido con la librería NumPy, de alrededor de 50 archivos \textit{.wav} de un segundo de duración.
Estos audios se generan sintéticamente
con frecuencias aleatorias y formas de onda básicas (\textit{sine}, \textit{square}, \textit{sawtooth}, \textit{triangle} y \textit{noise}).

En la fase de preprocesamiento y entrenamiento, desarrollada utilizando TensorFlow, se opta por
no aportar a la red formas de audio crudo. En su lugar, se extraen las características calculando
los Coeficientes Cepstrales en las Frecuencias de Mel (MFCC) de cada muestra. El modelo tiene el objetivo de clasificar 
y asociar la matriz de coeficientes con la etiqueta correspondiente a su tipo de onda. A pesar de esto, 
al plantear una arquitectura básica no convolucional, la red neuronal carece de los mecanismos para
realizar una interpretación correcta de los patrones bidimensionales de los MFCC. Como consecuencia, el 
modelo arroja una precisión muy baja. 

El \textbf{Prototipo 2} se basa en el Prototipo 1. Sigue atendiendo al objetivo de la clasificación, pero tiene 
características que difieren del anterior. Esta vez, los audios se convierten a espectrogramas \textit{.png}
mediante el uso de la librería de \textit{librosa}. De esta forma se consigue un menor coste computacional y de almacenamiento.
Al tratarse realmente de imágenes, se pueden utilizar modelos preentrenados de visión capaces de clasificarlas.

En el entrenamiento se hace uso de \textbf{ResNet34 preentrenada} vía \textit{fastAI}. Para su adaptación, se utilizan \textbf{DataBlocks},
una herramienta de \textit{fastAI} mediante la cual se especifican que las entradas son las imágenes (espectrogramas) y las salidas
son las categorías (tipo de onda). El etiquetado de la categoría de cada archivo se obtiene a partir su nombre y se divide
destinando un 80\% de las imágenes para entrenar y un 20\% para validación.

Como resultado de utilizar una red convolucional y la arquitectura externa \textbf{ResNet34}, el modelo obtuvo una buena precisión, especialmente 
identificando ondas puras dentro del rango entrenado. 

\subsection{Enfoque de Regresión}
La clasificación sirvió como prueba de concepto y como primer acercamiento a un modelo de procesamiento de audio, permitiendo que se explorasen tecnologías, biblotecas y métodos que serían de gran utilidad durante la totalidad de las fases del desarrollo. 

Sin embargo, para la implementación de un modelo que cumpliese con la descripción del trabajo, se decide dar un salto hacia la síntesis paramétrica continua. Para ello, fue necesario tomar una decisión
de diseño fundamental: la elección del motor de síntesis. Tras la evaluación de distintas alternativas, 
se optó por centrar el proyecto en la síntesis por Modulación de Frecuencia (FM). Esta técnica presenta
ventajas claras para el aprendizaje automático. A diferencia de la síntesis aditiva, que requiere decenas
de osciladores simultáneos, la síntesis FM es muy eficiente, consiguiendo generar timbres complejos
a partir de un número reducido de parámetros. De esta forma, se puede comenzar utilizando, por ejemplo, 
la frecuencia portadora, el ratio y el índice de modulación, y luego ampliar el número de variables para
conseguir sonidos más complejos. Estas características acotan la dimensionalidad del problema predictivo en un rango viable para la red neuronal.   

Si bien con la clasificación queda validada la extracción de características acústicas, una vez tomada 
esta última decisión, deja de ser útil. Para la inferencia de parámetros de la síntesis FM, resulta 
obligatoria la evolución hacia una arquitectura de regresión.

Para el \textbf{Prototipo 3} se genera un dataset mediano de aproximadamente 15000 muestras, construido 
mediante un barrido de parámetros en un sintetizador FM de \textbf{pyo}. Para la reducción de la carga computacional
y la optimización del espacio, se sustituye \textbf{librosa} por la librería de \textit{torchaudio}. Esta permite
transformar los archivos de audio (\textit{.wav}) en tensores de espectrogramas (\textit{.pt}),
asociados mediante el uso de un (\textit{.csv}) a sus valores numéricos. El 80\% del conjunto de datos se destina 
a entrenamiento y el otro 20\% a validación.

El entrenamiento se implementa mediante \textbf{PyTorch}. Se diseña una arquitectura convolucional 
propia (SmallCNNRegressor). La red neuronal se compone de dos bloques principales:
uno de extracción de características y una cabeza de regresión. 

El primer bloque recibe los espectrogramas en forma de tensores de un único canal (lo que equivaldría a 
imágenes en escala de grises) y los procesa en otros tres bloques convolucionales. Cada uno de estos aplica
kernels de tamaño 3x3 (o núcleos convolucionales) que se deslizan sobre la matriz del espectrograma
para extraer características espaciales, detectando así patrones visuales en la representación espectral del sonido.
A esto le sigue una función de activación no lineal \textbf{ReLU} y una capa de reducción espacial (\textbf{Max Pooling}). 
El último bloque convolucional termina en una capa de agrupamiento (\textbf{Adaptive Average Pooling}), que comprime los
mapas de características a una matriz 4x4.

El siguiente paso es la cabeza de regresión, en donde la información espacio-temporal del espectrograma
(extraída en el bloque anterior), se pasa a un vector de una única dimensión y atraviesa dos capas
completamente conectadas (\textbf{Perceptrón Multicapa}). La primera, comprime la información a 128 neuronas,
y la segunda, saca el vector por una dimensión de salida de tamaño 3, que se corresponde con los tres parámetros
del sintetizador: frecuencia portadora, ratio y el índice de modulación.

El modelo optimiza el proceso de los espectrogramas mediante el algoritmo Adam, utilizando el Error Cuadrático
Medio (MSELoss) como función de pérdida. Y, en cuanto al ciclo de entrenamiento, este se configuró con 
una duración de 5 épocas (\textit{epochs}), un número bajo pero suficiente para analizar el comportamiento del modelo. En este prototipo se decidió no hacer una búsqueda extensa
de hiperparámetros al tratarse de un prototipo intermedio, con un fin de experimentación inicial 
con la regresión.

A pesar de esto, se puede observar una precisión muy limitada en los resultados experimentales.
Se consigue mantener un MSE bajo en el \textit{Ratio} e \textit{Index}, pero el error en el
\textit{Carrier} resulta desproporcionado, debido a operar en la escala de miles de hercios. 
De este prototipo se puede sacar la conclusión de que minimizar el error matemático sobre los
parámetros no garantiza la similitud acústica del sonido generado y de la necesidad de trabajar sobre parámetros normalizados para contrarrestar las grandes diferencias en sus escalas. 


\section{El Problema de la Inyectividad y la Necesidad de una Arquitectura Compleja}
El Prototipo 3 contiene una limitación fundamental muy relevante en la sintesís FM:
su inyectividad es nula. Esto quiere decir que diferentes combinaciones de parámetros pueden dar lugar
a sonidos idénticos. Un ejemplo, radical pero que explica muy bien este resultado, es un caso en el que
el índice de modulación sea 0. En ese caso no es verdaderamente relevante el valor del ratio, sin embargo, 
la función de pérdida, implementada con \textbf{MSE}, penaliza gravemente que este difiera del original. 

Otro ejemplo, más común, ocurre en los casos en los que se trabaja sobre un índice de modulación alto, de tal forma que la frecuencia moduladora predomina sobre la portadora. Si tenemos un \textit{Carrier} de 1000 Hz y un \textit{Ratio} de 1.0, la frecuencia moduladora resultante será de 1000 Hz, generándose los armónicos en saltos de mil en mil hercios. Por otro lado, con un \textit{Carrier} de 2000 Hz y un \textit{Ratio} de 0.5, se obtiene la misma moduladora de 1000 Hz y, por lo tanto, los mismos armónicos. Todo esto sumado a un índice de modulación alto, harían que estos dos audios, paramétricamente muy distantes, suenen muy similares al oído humano. 
Este obstáculo hace que se tenga que replantear la función de pérdida.


\subsection{Prototipo 4: Red Convolucional de Doble Cabeza}
El \textbf{Prototipo 4} se diseña para solucionar el problema de la inyectividad, tratando de asegurar
que la inferencia de parámetros se corresponda con la fidelidad acústica. Este modelo se basa en una 
arquitectura convolucional de doble cabeza (\textbf{Encoder-Decoder}).

En cuanto al preprocesamiento del conjunto de datos, se implementa una mejora para evitar que la diferencia 
tan grande entre las escalas de difentes parámetros (como la frecuencia portadora frente al ratio) comprometa
el cálculo del gradiente. Los tensores se normalizan estadísticamente (\textbf{Z-score}), restando la media y 
dividiendo por la desviación estándar de cada parámetro. 

Para cada parámetro $x$, el valor normalizado $z$ se calcula según la Ecuación \ref{eq:zscore}:
\begin{equation}
z = \frac{x - \mu}{\sigma}
\label{eq:zscore}
\end{equation}
Donde $\mu$ representa la media aritmética de dicho parámetro en todo el dataset de entrenamiento y $\sigma$ su desviación estándar.

La topología de la red (CNNRegressor4), consta de un procesamiento dividido en varias fases:

Primero se pasa por un Codificador (\textbf{Encoder}), el cual consta de una secuencia de tres bloques
convolucionales que convierten el espectrograma inicial en una representación comprimida, extrayendo así sus 
características fundamentales. Para ello, cada bloque aplica núcleos (kernels) de 3x3. Posteriormente, 
pasa por una capa de normalización por lotes (Batch Normalization), que estabiliza y acelera el entrenamiento, y 
una activación no lineal ReLU. Finalmente, se aplica una reducción espacial en cada bloque (Max Pooling).
Con todo esto, la matriz se ha reducido a una octava parte de su tamaño original y la profundidad de la red
ha incrementado de 1 a 128 canales. Como resultado, se obtiene una representación comprimida
rica en características acústicas.

Al salir del \textbf{Encoder}, el flujo de datos atraviesa lo que se denomina como
cuello de botella (\textbf{Bottleneck}). Este bloque convolucional no
aplica una reducción espacial, manteniendo estables las dimensiones de la matriz. Su función es duplicar el número
de canales (profundidad del tensor), pasando de 128 a 256 canales. Esto se hace con la finalidad de 
mezclar los patrones que el \textbf{Encoder} ha extraído previamente, creando una representación
densa y unificada.

Después se pasa por la fase de Inferencia Numérica, donde ocurre la regresión de los parámetros del sintetizador FM
(frecuencia portadora, ratio e índice). Se aplica una capa de agrupamiento global (Adaptive Average Pooling) ya que los 
parámetros describen el sonido completo y no un punto temporal concreto en el espectrograma.
Para ello, se calcula la media de todas las posiciones de cada canal, quedándose con un único
número como resumen por canal. Finalmente, el vector unidimensional, resultado del proceso anterior,
atraviesa un \textbf{Perceptrón Multicapa}, que lo proyecta hacia las tres neuronas de salida.

Por último, se encuentra la fase de Reconstrucción (\textbf{Decoder}), la cual se encarga de reconstruir 
el espectrograma original a partir de la representación comprimida del \textbf{Bottleneck}. 
Como el \textbf{Encoder} comprime tres veces, el \textbf{Decoder} descomprime tres veces. Para ello,
utiliza tres capas convolucionales transpuestas (ConvTranspose2d).

El motivo para utilizar el \textbf{Decoder} es como método de regularización forzada del \textbf{Encoder}.
En el caso en el que el codificador borre información importante, como armónicos o patrones temporales, que no afecte directamente a los parámetros
de la regresión, el (\textbf{Decoder}) no será capaz de redibujar correctamente el espectrograma, disparando así
el error en la función de pérdida híbrida. Aparte, si el \textbf{Bottleneck} no guarda suficiente información, ocurrirá lo mismo.
Todo esto fuerza, mediante el cálculo de gradientes, a que las capas iniciales aprendan a conservar el timbre y la representación
del sonido.

\subsection{Prototipo 4: Función de Pérdida Híbrida}
Con el objetivo de solucionar el problema de la inyectividad, se consolida el uso de utilizando
Función de Pérdida Híbrida, abandonando el Error Cuadrático Medio (MSE). Esta función evalúa la pérdida
de dos formas simultáneamente, dándole prioridad a la similitud y convergencia de espectrogramas (pérdida y convergencia espectral)
y dejando en segundo plano la similitud de los parámetros (pérdida paramétrica).

Por un lado tenemos la \textbf{Pérdida y Convergencia Espectral}. Estas funciones son las principales
para la composición del error, teniendo mucho más peso que la pérdida paramétrica, y estando encargadas de asegurar
que un sonido generado sea acústicamente similar al original. 

La \textbf{pérdida espectral} se calcula mediante el Error Cuadrático Medio (\textbf{L1 Loss}), el cual mide la diferencia directa entre 
espectrogramas (píxel a píxel) en escala logarítmica (dB), siendo la métrica de mayor peso para el error. Con esto se garantiza unos picos de energía y
formas generales que coincidan en tiempo y frecuencia.

Matemáticamente, si $S$ es el espectrograma original en decibelios y $\hat{S}$ el reconstruido, ambos con $N$ elementos, la pérdida L1 se define como:
\begin{equation}
L_1 = \frac{1}{N} \sum_{i=1}^{N} |S_i - \hat{S}_i|
\label{eq:l1_loss}
\end{equation}

En cuanto a la \textbf{convergencia espectral}, es un refuerzo (tiene un peso bajo) para la pérdida espectral, que compara las magnitudes de forma relativa con la norma de
Frobenius, capturando la estrutura global del espectro (como los armónicos y la energía). Estas características 
son muy relevantes para encontrar una representación fiel al timbre original. 

La convergencia espectral ($C_S$) emplea la norma de Frobenius ($||\cdot||_F$) sobre las magnitudes matriciales para medir la distancia relativa entre las energías, tal y como detalla la Ecuación \ref{eq:spectral_convergence}:
\begin{equation}
C_S = \frac{||S - \hat{S}||_F}{||S||_F}
\label{eq:spectral_convergence}
\end{equation}

Por último se encuentra la \textbf{Pérdida Paramétrica}, aplicada mediante \textbf{SmoothL1}. Esta penaliza las diferencias en los parámetros
del sintetizador, por lo que no se elimina totalmente la regresión numérica de los parámetros, pero se le da un peso
considerablemente bajo. 

La Ecuación \ref{eq:smooth_l1} muestra el comportamiento por tramos de SmoothL1 para una predicción $\hat{y}$ frente a su valor real $y$:
\begin{equation}
L_{smooth} = 
\begin{cases} 
0.5 (y - \hat{y})^2, & \text{si } |y - \hat{y}| < 1 \\ 
|y - \hat{y}| - 0.5, & \text{en caso contrario} 
\end{cases}
\label{eq:smooth_l1}
\end{equation}

La ponderación de los pesos en la función de pérdida híbrida es fundamental para el correcto aprendizaje del Modelo
y también para solucionar el problema de la inyectividad. Se le otorga un peso completo (1.0) a la pérdida espectral,
un peso medio (0.5) a la convergencia espectral y un peso muy bajo (0.05) a la pérdida paramétrica. De esta manera, en 
el caso de que la red infiera una serie de parámetros que difieran considerablemente de los parámetros originales, si sus 
espectrogramas son muy similares (y, por tanto, su tiembre también lo es), el error será mínimo. El modelo debe memorizar números
exactos, si no que analiza la imagen del sonido, su espectrograma.